% ============================================================
%  chapters/ch09-action-histories.tex
% ============================================================

\chapter{Action Histories and Spherepop Dynamics}

\chapterepigraph{%
  The present moment is always the consequence\\
  of everything that has already been decided.\\
  Action history is the accumulation of constraint.
}{Flyxion, \textit{Semantic Infrastructure}}

\sidenote{Connects to\\decision theory,\\Lagrangian mechanics,\\and RSVP action\\(App.\,K §K.20)}

The previous four chapters developed Spherepop as a static and kinematic
system: expressions, their syntax, their evaluation, their geometry, their
renormalization. This chapter introduces a dynamic structure that goes deeper:
the \emph{action history} of a computation, and the variational principle
that governs it.

\section{The Spherepop Lagrangian}

Classical mechanics describes particle trajectories as the extrema of
an action functional:
\[
  \mathcal{S}[\gamma] = \int_0^T L(q, \dot{q})\,dt.
\]
The Euler–Lagrange equations then determine which trajectories are physical.

Spherepop evaluation has an analogous structure. Rather than a particle
moving through physical space, we have an expression evolving through
expression space. Rather than kinetic energy minus potential energy,
we have evaluation cost minus structural richness.

\begin{definition}{Spherepop Lagrangian}{sp-lagrangian}
  For a reduction sequence $\sigma = (E_0 \to E_1 \to \cdots \to E_n)$,
  the \emph{Spherepop Lagrangian} is:
  \[
    L(E_t, \dot{E}_t) = C(E_t) - \lambda \cdot \mathrm{Acc}(E_t)
  \]
  where $C(E_t)$ is the complexity (number of non-atomic spheres),
  $\mathrm{Acc}(E_t) = \log |\mathcal{A}(E_t)|$ is the accessibility
  entropy, and $\lambda > 0$ is a tradeoff parameter.
\end{definition}

The action of the reduction sequence is:
\[
  \mathcal{S}[\sigma] = \sum_{t=0}^{n-1} L(E_t, \dot{E}_t).
\]
Evaluation minimizes this action: among all sequences reaching the
terminal expression, evaluation by the innermost-first (depth-first)
strategy minimizes $\mathcal{S}$.

\section{Action Histories}

The concept of an action history extends the variational picture.
A computation does not exist only in the present moment — it carries the
entire record of how it arrived at its current state.

\begin{definition}{Action History}{action-history}
  The \emph{action history} of an expression $E_t$ in a reduction
  sequence is the sequence of all previous states and pop operations:
  \[
    \mathcal{H}_t = (E_0, p_0, E_1, p_1, \ldots, p_{t-1}, E_t)
  \]
  where $p_i$ denotes the pop operation at step $i$.
\end{definition}

The action history is the computational analog of the path integral
in quantum mechanics. Just as the quantum amplitude for a particle to
reach state $x$ at time $T$ is a sum over all paths from the initial
state to $x$, the ``weight'' of a computation reaching expression $E_t$
is determined by the action of all paths leading to it.

\section{Commitment Variables}

\begin{remark}
  \textbf{The $0.73\,$Hz loop and action history.}
  Cognitive neuroscience identifies a characteristic frequency near
  $0.73\,$Hz in the aggregation of sensory inputs into conscious
  perception — the rate at which the brain integrates distributed
  neural activity into a unified temporal snapshot of the present moment.
  In the action-history framework, this loop is the rate of
  \emph{commitment}: each cycle converts a portion of remaining freedom
  $F_t$ into committed past $\Gamma_t$. The $0.73\,$Hz rhythm is the
  cognitive analogue of the cosmic entropy clock — the rate at which
  the agent's future collapses into its past, one irreversible snapshot
  at a time. An agent's ``present moment'' is the temporal resolution
  of its commitment cycle, not a physical primitive.
\end{remark}


A key concept in action history theory is the \emph{commitment variable}:
a choice made at some point in the evaluation that is now fixed and
cannot be undone.

\begin{definition}{Commitment Variable}{commitment-var}
  A \emph{commitment} is a pop step $p_i$ in an action history
  $\mathcal{H}_t$ (with $i < t$). Once committed, the popped sphere
  is gone — the information it contained has been discarded. The
  commitment is irreversible.
  The \emph{commitment set} $\Gamma_t = \{p_0, p_1, \ldots, p_{t-1}\}$
  is all decisions made up to time $t$.
\end{definition}

The commitment set accumulates throughout evaluation. Early commitments
(deep pops) constrain the behavior of later, shallower pops. This
is why evaluation order matters for efficiency even when it does not
matter for correctness: different evaluation orders make different
commitments at different times, and some commitment orders allow
subsequent decisions to be made more efficiently.

\section{Remaining Freedom}

The complement of the commitment set is the \emph{remaining freedom}:
the set of choices not yet made.

\begin{definition}{Remaining Freedom}{remaining-freedom}
  The \emph{remaining freedom} at time $t$ in an evaluation sequence
  is the set of expressions reachable from $E_t$ by some evaluation
  strategy:
  \[
    \mathcal{F}_t = \{E' : E_t \to_\text{pop}^* E'\}.
  \]
  The \emph{freedom measure} is $F_t = |\mathcal{F}_t|$.
\end{definition}

\begin{proposition}{Freedom Decreases with Commitment}{freedom-decrease}
  For any evaluation sequence, $F_{t+1} \leq F_t$.
\end{proposition}

\begin{proof}
  $\mathcal{F}_{t+1} \subseteq \mathcal{F}_t$: every expression reachable
  from $E_{t+1}$ is also reachable from $E_t$ (via the additional pop
  step $E_t \to E_{t+1}$). But $\mathcal{F}_{t+1}$ may exclude expressions
  reachable from $E_t$ by a different pop step, so the containment
  may be strict.
\end{proof}

This proposition formalizes a profound principle: \emph{computation
consumes freedom}. Each pop step irrevocably narrows the future.
This is the computational analog of the thermodynamic increase of
entropy: the past grows larger; the future shrinks.

\section{Accessibility Measures}

The accessibility measures of Chapter~1 and Appendix~J apply directly
to evaluation sequences. The accessibility entropy of an expression
is the logarithm of its freedom measure:
\[
  S_t = \log F_t = \log |\mathcal{F}_t|.
\]

\begin{proposition}{Accessibility Entropy Decreases Under Evaluation}{access-decrease}
  For any evaluation sequence, $S_{t+1} \leq S_t$.
\end{proposition}

This is the computational Second Law: evaluation is an irreversible
process that consumes accessibility. The initial expression has maximal
accessibility (many possible futures); the terminal atom has zero
accessibility (no future — computation is complete).

\begin{figure}[h]
  \centering
  \begin{tikzpicture}[scale=0.85]
    % Axes
    \draw[->, thick, RSVPdeep] (0,0) -- (6.5,0) node[right] {$t$ (steps)};
    \draw[->, thick, RSVPdeep] (0,0) -- (0,4.5) node[above] {measure};
    % C(E_t) — complexity, decreasing
    \draw[RSVPaccent, very thick]
      (0,4.0) -- (1,4.0) -- (1,3.2) -- (2,3.2) -- (2,2.3) --
      (3,2.3) -- (3,1.4) -- (4,1.4) -- (4,0.5) -- (5,0.5) -- (5,0);
    \node[RSVPaccent, font=\small\itshape] at (5.8, 2.0) {$C(E_t)$};
    % S_t — accessibility entropy, also decreasing
    \draw[RSVPmid, very thick, dashed]
      (0,3.5) to[out=0,in=170] (2,2.6) to[out=-10,in=165] (4,1.0)
      to[out=-15,in=175] (5,0);
    \node[RSVPmid, font=\small\itshape] at (5.8, 1.2) {$S_t$};
    % Labels
    \foreach \x in {0,1,2,3,4,5}
      \node[below, font=\tiny] at (\x,0) {$\x$};
    \node[left, font=\tiny] at (0,4.0) {$n$};
    \node[left, font=\tiny] at (0,0) {$0$};
    \node[RSVPdeep, font=\footnotesize] at (2.5,-0.6) {pop step};
  \end{tikzpicture}
  \caption{Complexity $C(E_t)$ and accessibility entropy $S_t$ as
    functions of evaluation step $t$. Both decrease monotonically
    to zero at the terminal state. $C$ decreases by exactly 1 per
    step; $S$ decreases more smoothly, capturing the branching
    structure of the evaluation space.}
  \label{fig:entropy-decrease}
\end{figure}

\section{Connection to Decision Theory}

The action history formalism connects directly to decision theory under
uncertainty. Consider a computation where some operators are stochastic:
they produce different outputs with different probabilities. The evaluation
then defines a probability tree, and the action history becomes a sample
path through this tree.

\begin{definition}{Stochastic Evaluation}{stoch-eval}
  A \emph{stochastic Spherepop system} replaces deterministic operators
  with probability distributions over outputs. The \emph{expected action}
  of an evaluation strategy $\pi$ is:
  \[
    \mathbb{E}_\pi[\mathcal{S}] = \sum_\sigma P_\pi(\sigma) \cdot \mathcal{S}[\sigma]
  \]
  where the sum is over all possible reduction sequences $\sigma$.
\end{definition}

Optimal decision-making — choosing the evaluation strategy $\pi^*$ that
minimizes expected action subject to accessibility constraints — is
equivalent to solving a Markov decision process defined on the expression
manifold.

\section{Connection to RSVP}

The action history formalism is the bridge between Spherepop and RSVP.
In RSVP, the action functional is:
\[
  \mathcal{S}_\text{RSVP}[\Phi,\mathbf{v},S] =
  \int \mathcal{L}(\Phi,\mathbf{v},S)\,d^4x
\]
(Appendix~K, Definition~K.20). The Spherepop Lagrangian $L(E_t)$ is
the discrete analog of $\mathcal{L}$:

\begin{itemize}
  \item Complexity $C(E_t)$ $\leftrightarrow$ scalar density $\Phi$:
        both measure the concentration of unresolved structure.
  \item Pop operations $\leftrightarrow$ vector flow $\mathbf{v}$:
        both describe directed motion through the state space.
  \item Accessibility entropy $S_t$ $\leftrightarrow$ entropy field $S$:
        both measure the volume of available futures.
\end{itemize}

The RSVP field equations (Appendix~K) are the continuous limit of the
Spherepop action history: as the expression space becomes continuous
and evaluations become infinitesimally small, the discrete action sum
$\sum L(E_t)$ becomes the continuous action integral $\int \mathcal{L}\,d^4x$.

\begin{namedthm}{Spherepop–RSVP Correspondence}
  The Spherepop action history formalism is the discrete skeleton of
  the RSVP field theory. Expressions correspond to field configurations;
  evaluation sequences correspond to field trajectories; complexity
  corresponds to scalar density; accessibility corresponds to entropy.
  The two frameworks are related by a limiting process: RSVP is
  the continuum limit of Spherepop dynamics.
\end{namedthm}

\begin{exercise}
  Compute the action $\mathcal{S}[\sigma]$ for the evaluation
  of $1 + (3 \times (2 \uparrow 2))$, using the Lagrangian
  $L(E_t) = C(E_t) - S_t$ with $\lambda = 1$. Assume $S_0 = \log 6$,
  $S_1 = \log 5$, $S_2 = \log 3$, $S_3 = 0$.
\end{exercise}

\begin{exercise}
  The freedom measure $F_t = |\mathcal{F}_t|$ measures the number of
  terminal expressions reachable from $E_t$. For a deterministic,
  confluent operator set, show that $F_t = 1$ for all $t \geq 0$.
  (What does this say about accessibility in fully deterministic
  computation?)
\end{exercise}

\begin{exercise}[title={Open problem}]
  Define an appropriate notion of \emph{Spherepop entropy production}:
  the rate at which accessibility is consumed per pop step. For which
  evaluation strategies is entropy production minimal? Does minimal
  entropy production correspond to any known computational optimization
  strategy?
\end{exercise}
